\chapter{PENDAHULUAN}

\section{Latar Belakang}

% TODO: Tulis latar belakang penelitian Anda di sini.
% Jelaskan konteks, urgensi, dan relevansi topik yang dipilih.
% Rujuk pada penelitian terdahulu dan data sekunder jika ada.

Sistem informasi manajemen asuransi kesehatan sumber terbuka seperti openIMIS (\textit{open Insurance Management Information System}) dirancang untuk mendukung proses administratif penjaminan kesehatan, misalnya pendaftaran peserta, penelusuran informasi peserta, dan pengelolaan klaim. openIMIS merupakan perangkat lunak yang telah disertifikasi sebagai digital public good dan dilaporkan digunakan di 14 negara untuk mengelola program yang menjangkau lebih dari 39 juta penerima manfaat, mencakup asuransi kesehatan, asuransi kecelakaan kerja, transfer tunai, dan skema voucer \citet{openimis2026}. Antarmuka sistem semacam ini memuat elemen informasi utama yang menopang pelaksanaan proses tersebut, seperti informasi eligibilitas peserta dan status klaim. Apabila elemen-elemen tersebut tidak berada pada posisi yang secara alami menarik atensi visual pengguna, pelaksanaan proses administratif berpotensi kurang efisien atau kurang tepat, namun sejauh ini belum ada mekanisme sistematis untuk memeriksa apakah elemen informasi utama memang berada pada zona yang diprediksi menarik atensi visual.

Isu ketidaksesuaian antara desain antarmuka dan kebutuhan pengguna dalam sistem informasi kesehatan bukan hal baru. Kajian terhadap implementasi OpenMRS selama 15 tahun di negara berpendapatan rendah dan menengah menemukan bahwa hambatan adopsi sistem rekam medis elektronik tidak hanya perihal infrastruktur, melainkan juga \textit{usability} antarmuka dan rendahnya literasi teknologi informasi tenaga kesehatan, sehingga memerlukan pendekatan \textit{Human-Centered Design} \citep{verma2021openmrs}. Selain itu, kajian pada IMIA Yearbook menegaskan bahwa \textit{resilience, accessibility,} dan \textit{equity} menjadi tiga pilar utama sistem kesehatan digital, dan menyoroti upaya pembangunan ulang arsitektur antarmuka OpenMRS generasi baru dengan pengujian pengguna yang mencakup lintas negara \citep{paton2022open}. Kedua kajian ini menunjukkan bahwa evaluasi \textit{usability} pada sistem informasi kesehatan penting, namun keduanya masih bertumpu pada metode evaluasi konvensional (studi lapangan, pengujian pengguna langsung) yang memerlukan sumber daya besar sehingga sulit direplikasi dalam skala kecil seperti dalam konteks penelitian tugas akhir.

Di sisi lain, bidang prediksi atensi visual (\textit{visual attention/gaze saliency prediction}) berbasis \textit{deep learning} telah berkembang matang selama kurang lebih satu dekade, dengan sejumlah model seperti DeepGaze, UMSI, dan TranSalNet, serta didukung oleh dataset benchmark antarmuka digital berskala besar seperti UEyes yang dipublikasikan pada ACM CHI 2023. Model-model ini memungkinkan estimasi peta atensi visual pada sebuah antarmuka tanpa memerlukan partisipan manusia di setiap sesi pengujian, berbeda dengan pendekatan berbasis riil (misalnya \textit{heatmap} berbasis klik/\textit{mouse-tracking}) yang selama ini lebih umum digunakan dalam evaluasi usabilitas. Karakteristik inilah yang menjadikan model prediksi atensi visual berpotensi sebagai pendekatan analisis yang lebih efisien dari sisi biaya dan waktu dibandingkan studi eye-tracking konvensional.

Meskipun demikian, penelusuran literatur pada Scopus dan IEEE Xplore, penelusuran sitasi lanjut terhadap UEyes pada Scopus dan Google Scholar, serta pencarian web eksploratori pada Oktober 2026 (selanjutnya disebut literatur yang ditelusuri) belum menemukan studi yang menerapkan model prediksi atensi visual pada antarmuka sistem informasi kesehatan administratif, termasuk pada openIMIS. Studi pada domain kesehatan dan asuransi yang ditemukan umumnya memakai pelacakan mata secara manual atau menganalisis citra diagnostik.Kekosongan ini konsisten dengan kendala praktis di lapangan: tidak tersedia dataset eye-tracking atau data interaksi pengguna openIMIS yang dapat dijadikan basis pelatihan atau validasi kuantitatif standar (metrik seperti NSS, CC, atau KLD), dan pengumpulan data semacam itu secara mandiri berada di luar kelayakan waktu dan sumber daya penelitian tugas akhir tingkat sarjana. Akses ke pengguna operasional maupun pakar domain untuk keperluan validasi juga sulit diperoleh.

Untuk menjawab kendala tersebut, penelitian ini mengusulkan pendekatan yang tidak bergantung pada \textit{fine-tuning} berbasis data domain-spesifik. Model prediksi atensi visual pra-latih (\textit{pretrained}), yaitu TranSalNet\_Res, digunakan langsung tanpa modifikasi bobot. Kredibilitas umum model diuji terlebih dahulu pada benchmark publik UEyes, khususnya kategori \textit{Webpage} dan \textit{Desktop UI} yang secara struktural paling mendekati antarmuka openIMIS, sebagai tolok ukur independen, bukan sebagai basis pelatihan untuk domain openIMIS. Model yang sama kemudian diterapkan pada tangkapan layar openIMIS untuk menghasilkan peta atensi. Skor saliency prediktif pada area setiap elemen antarmuka, yang kotaknya diambil dari struktur halaman (\textit{Document Object Model}/DOM), selanjutnya dibandingkan antara elemen informasi utama, yang dilabeli secara sistematis dari dokumentasi resmi, materi pelatihan, dan video demonstrasi openIMIS menggunakan buku kode yang terdokumentasi, dan elemen yang tidak berlabel. Perlu ditegaskan bahwa model yang digunakan memprediksi atensi visual yang bersifat \textit{bottom-up} pada kondisi pengamatan bebas (\textit{free-viewing}), sebagaimana data acuan UEyes yang dikumpulkan tanpa tugas spesifik (Jiang et al., 2023), sehingga hasil penelitian ini hanya menyatakan tingkat menonjolnya elemen secara visual (saliency) menurut prediksi model dan tidak mengklaim perilaku pengguna yang sedang mengerjakan tugas. Topik penelitian ini memiliki relevansi kuat dengan roadmap riset \textit{Human-Computer Interaction} dan penerapan kecerdasan buatan pada sistem informasi kesehatan di Departemen Ilmu Komputer dan Elektronika FMIPA UGM, khususnya arah riset yang dikembangkan pada Lab Rekayasa Perangkat Lunak dan Data. Riset ini diharapkan dapat memberikan kontribusi metodologis maupun praktis terhadap diagnosis risiko desain antarmuka pada sistem informasi kesehatan administratif.


\section{Rumusan Masalah}

% TODO: Tulis rumusan masalah dalam bentuk pertanyaan penelitian.
% Rumusan masalah harus spesifik, terukur, dan dapat dijawab.

Berdasarkan latar belakang yang telah diuraikan, dapat dirumuskan permasalahan penelitian sebagai berikut:
\begin{enumerate}
  \item Antarmuka sistem administrasi kesehatan seperti openIMIS memuat elemen informasi utama yang menopang pelaksanaan proses administratif, namun tingkat menonjolnya elemen tersebut secara visual dalam peta atensi visual prediktif perlu dievaluasi secara sistematis.
  \item Model \textit{deep learning} untuk prediksi atensi visual telah berkembang matang pada berbagai domain antarmuka digital umum, namun belum ditemukan acuan mengenai penerapannya pada antarmuka sistem administrasi kesehatan berbasis asuransi pada literatur yang ditelusuri, sehingga tingkat kredibilitas model pra-latih TranSalNet\_Res pada benchmark publik UEyes kategori \textit{Webpage} dan \textit{Desktop UI} perlu diukur sebagai dasar menafsirkan peta atensi pada antarmuka openIMIS.
  \item Ketidaktersediaan dataset \textit{eye-tracking} atau data interaksi pengguna yang spesifik pada domain openIMIS menjadi kendala untuk validasi kuantitatif standar (NSS/CC/KLD) untuk model prediksi atensi visual, sehingga diperlukan pendekatan metodologis alternatif yang transparan dan dapat direplikasi untuk membandingkan atensi visual prediktif dengan elemen informasi utama, tanpa memerlukan fine-tuning berbasis data domain dan tanpa validasi pakar.
\end{enumerate}

\section{Batasan Masalah}

% TODO: Tulis batasan masalah untuk memfokuskan penelitian.
% Batasan dapat mencakup ruang lingkup, teknologi, atau aspek yang tidak diteliti.

Untuk menjaga fokus dan kelayakan penelitian dalam lingkup tugas akhir, ditetapkan batasan sebagai berikut:

\begin{enumerate}
  \item Objek studi dibatasi pada antarmuka sistem openIMIS versi \textit{sandbox}/demo, bukan lingkungan produksi aktif
  \item Model prediksi atensi visual yang digunakan adalah TranSalNet\_Res pra-latih \textit{open-source} tanpa \textit{fine-tuning} tambahan menggunakan data domain openIMIS. Model lain (mis. TranSalNet\_Dense, SUM) dan pendekatan zero-shot berbasis model visual-language berada di luar cakupan inti dan hanya disebut sebagai pengembangan lanjutan.
  \item Validasi akurasi umum model dilakukan dengan publik UEyes, khusus kategori \textit{Webpage} dan \textit{Desktop UI}, sebagai ukuran kredibilitas umum dan bukan sebagai metrik akurasi domain spesifik openIMIS
  \item Elemen informasi utama pada antarmuka openIMIS ditentukan melalui pelabelan oleh satu pelabel (peneliti) berdasarkan buku kode yang terdokumentasi, dengan sumber berupa dokumentasi resmi, materi pelatihan, dan video demonstrasi resmi openIMIS. Suatu elemen dilabeli apabila didukung oleh minimal dua sumber. Penelitian ini tidak melibatkan validasi oleh pakar atau pengguna operasional, sehingga hasilnya bersifat eksploratif dan tidak mengklaim kesesuaian dengan perilaku pengguna nyata.
  \item Sumber pelabelan dan tangkapan layar dapat berasal dari versi antarmuka yang berbeda. Pemetaan elemen dilakukan berdasarkan fungsinya, dan perbedaan versi dinyatakan sebagai keterbatasan penelitian.
  \item Luaran penelitian berupa analisis yang dituangkan dalam laporan skripsi; penelitian tidak mencakup pengembangan aplikasi atau perangkat lunak bagi pengguna akhir. Inferensi model, perhitungan skor saliency per elemen, dan analisis statistik dikerjakan melalui skrip analisis (\textit{notebook}) yang didokumentasikan sebagai instrumen penelitian. Pengembangan backend/infrastruktur openIMIS itu sendiri berada di luar cakupan penelitian.
  \item Penelitian tidak mencakup siklus \textit{redesign} penuh terhadap antarmuka openIMIS, keluaran penelitian adalah analisis tingkat menonjolnya elemen secara visual dan kandidat rekomendasi, bukan implementasi ulang desain.
  \item Validasi kuantitat skala kecil dengan partisipan proksi (5-10 orang) merupakan opsi penguat apabila logistik memungkinkan, bukan komponen wajib penelitian.
  \item Cakupan layar dibatasi pada sejumlah alur tugas yang dipilih berdasarkan kriteria kecukupan sumber (dijelaskan pada Bab 3), dan analisis hanya mencakup elemen pada viewport awal setiap layar.
  \item Klaim penelitian dibatasi pada atensi visual prediktif (bottom-up, free-viewing); penelitian tidak mengklaim perilaku atau atensi pengguna yang sedang mengerjakan tugas.
\end{enumerate}

\section{Tujuan Penelitian}

% TODO: Tulis tujuan penelitian yang ingin dicapai.
% Tujuan harus sesuai dengan rumusan masalah.

Berdasarkan rumusan masalah dan batasan yang telah diidentifikasi, tujuan penelitian ini adalah:

\begin{enumerate}
  \item Mengukur kredibilitas umum model prediksi atensi visual yang dipilih melalui benchmark publik UEyes pada kategori \textit{Webpage} dan \textit{Desktop UI}
  \item Membandingkan tingkat menonjolnya secara visual (berdasarkan skor saliency prediktif) elemen informasi utama dengan elemen tidak berlabel pada antarmuka openIMIS, berdasarkan pelabelan dari dokumentasi resmi, materi pelatihan, dan video demonstrasi, menggunakan ukuran persentil dan uji permutasi dengan hipotesis dua arah.
  \item Merumuskan temuan beserta kandidat rekomendasi perbaikan tata letak antarmuka yang memerlukan pengujian lanjutan dengan pengguna, dengan menyatakan keterbatasan validasinya secara eksplisit.
\end{enumerate}


\section{Manfaat Penelitian}

% TODO: Tulis manfaat penelitian bagi akademik dan praktis.

Penelitian ini diharapkan memberikan manfaat sebagai berikut:

\begin{enumerate}
  \item \textbf{Bagi akademik:} memberikan kontribusi metodologis awal pada irisan bidang prediksi atensi visual dan desain antarmuka sistem informasi kesehatan administratif, sebuah gap yang belum ditemukan pada literatur yang ditelusuri.
  \item \textbf{Bagi pengembang/pengelola openIMIS:} menyediakan pendekatan penelusuran awal yang efisien dari sisi biaya dan waktu untuk mengidentifikasi kandidat risiko desain antarmuka yang perlu diuji lebih lanjut dengan pengguna, tanpa memerlukan pengumpulan data eye-tracking yang berbiaya tinggi dan berdurasi panjang, dan bukan sebagai pengganti uji pengguna.
  \item \textbf{Bagi komunitas pengembang dan pengelola sistem informasi asuransi kesehatan \textit{open source}:}Menjadi dasar awal bagi kajian usability lanjutan yang lebih besar dan berbasis data domain-spesifik di pengembangan selanjutnya.
  \item \textbf{Bagi peneliti selanjutnya:} memberikan \textit{baseline} metodologi (validasi \textit{benchmark} umum dan expert judgment tanpa \textit{fine-tuning}) yang berpotensi direplikasi pada domain sistem informasi lain dengan keterbatsan data serupa. 
\end{enumerate}

